Detecting unusual financial behaviour by comparing a transaction with a person's own history raises a question that deployed systems settle by fiat: how many transactions make a personal baseline worth using? We measure it. Data-efficiency curves for six unsupervised detectors, evaluated on 300 real bank accounts, a 240-person controlled simulation and 300 entities of a synthetic card dataset, show that a shrinkage-regularised Mahalanobis detector reaches 90\% of its best ranking performance (precision-recall AUC 0.657) at 25 personal transactions and gains nothing measurable beyond 100; the threshold is identical across five independent cohort draws. Labels come from provenance, not from any detector: a held-out transaction is replaced by a real transaction from a different account. Below about 10 transactions a population model outperforms personalisation, and a history-weighted shrinkage hybrid removes that penalty at no measurable cost later. Operating-point behaviour depends on how the alert threshold is calibrated: scoring the model's own training data makes the detector flag 96\% of transactions at 10 personal transactions, whereas leave-one-out calibration leaves it nearly silent there (1.2\%) and close to its nominal false-positive rate from 25. Detector rankings under scarcity reverse between datasets, and behavioural variability predicts the requirement under controlled manipulation (Spearman correlation 0.48 to 0.71) but not across real accounts. An anomaly here is a deviation from a baseline, not fraud.
